\section{Finite-time area transfer}\label{sec:area-transfer}

The two deformations used below can contract a metric on a set of small
weighted volume.  The following argument shows that this contraction
cannot substantially decrease the minimum enclosing area.  Its constants
use a scalar-curvature lower bound and the topology of the original
exterior.  No curvature bound for the changing metrics, and no limit of a
hyperbolic Hawking mass, is needed.

\subsection{Caps, intrinsic boundaries, and the initial surface}

Attach a compact orientable handlebody to each component of the original
boundary $S$.  Write $O$ for the union of these caps and $X$ for the
resulting connected manifold without boundary.  This operation is
topological: whenever a comparison metric is given, extend it smoothly
across $S$ into the caps.  The end is unchanged, and the extension is
complete.  The topology of $X$ and $O$ will remain fixed.  In particular,
\begin{equation}\label{eq:flow-topology-constant}
 b_2=\dim_{\mathbb F_2}H_2(X\setminus\operatorname{int}O;\mathbb F_2)<\infty.
\end{equation}
Finiteness follows by retracting the product end onto its inner sphere;
what remains is a compact manifold with boundary.

We will start from a largest perimeter-minimizing enclosure $E_0$ of
$O$ alone.  In each application, the obstacle and end barriers establish
that its boundary is smooth, free, minimal, and disjoint from $O$.
Thus $O\subset E_0$ and $\partial O\subset\operatorname{int}E_0$.
The open representative $E_0$ is precompact.  Each component meets $O$,
since an extra component could be removed with a strict decrease of
perimeter.  Bounded components of its complement can similarly be
filled.  A largest minimizer is strictly outward minimizing: a larger
set with equal perimeter would also be a minimizer, contrary to
maximality.  This assertion concerns all precompact competitors, because
the far-out barriers allow any such competitor to be trimmed into the
fixed truncation used in constructing $E_0$.

The cap construction respects the intrinsic-boundary convention in
Section~\ref{sec:hypotheses}.  If $D$ is any admissible closed end-side
domain in the original exterior, then
$X\setminus\operatorname{int}_X D$ is a precompact finite-perimeter
enclosure of $O$.  Its perimeter is bounded by the area of the
\emph{entire} intrinsic boundary of $D$.  In particular, portions along
$S$ are counted; for $D=\Omega$ the perimeter is $|S|$, not zero.
Conversely, a smooth enclosure of $O$ whose boundary is disjoint from
$O$ and whose exterior is connected gives an admissible cut by taking
the closure of its exterior.  These observations also apply when the
surfaces or the caps have several components.

\begin{lemma}[Adjustment at the initial surface]\label{lem:flow-initial-adjustment}
Let $\widehat h$ be a smooth complete metric on $X$, and suppose $E_0$
has the properties just described.  If
$\Scal_{\widehat h}\ge-C_R$ on $X\setminus E_0$, then an arbitrarily
small, compactly supported smooth conformal change gives a metric
$\widetilde h$ for which
\begin{equation}\label{eq:flow-initial-adjustment}
 |\partial E_0|_{\widetilde h}=|\partial E_0|_{\widehat h},\qquad
 H_{\widetilde h}>0,\qquad
 \int_{\partial E_0}H_{\widetilde h}^{2}\,\dd A_{\widetilde h}\le1,
 \qquad \Scal_{\widetilde h}\ge-C_R-1
 \quad\hbox{on }X\setminus E_0.
\end{equation}
The set $E_0$ remains strictly outward minimizing.  The size of the
change may be chosen separately for every comparison metric.
\end{lemma}

\begin{proof}
In a two-sided signed-distance collar of $\partial E_0$, let $q$ be
positive on the exterior and choose
$\chi=q\zeta(q)$, where $\zeta$ is a nonnegative smooth cutoff equal to
one near zero.  Extend $\chi$ by zero beyond the collar.  Then
$\chi=0$ and $\partial_\nu\chi=1$ at the initial surface, while
$\chi\ge0$ outside $E_0$.  Set
$\widetilde h=e^{2a\chi}\widehat h$, with $a>0$.  The conformal
mean-curvature formula gives $H_{\widetilde h}=2a$ on $\partial E_0$.
Consequently its area is unchanged and its squared-mean-curvature
integral is $4a^2|\partial E_0|_{\widehat h}$.

Every outward competitor has its boundary outside $E_0$, where the
metric has only increased, so strict outward minimality is preserved.
Finally,
\begin{equation}\label{eq:flow-conformal-scalar}
 \Scal_{\widetilde h}
 =e^{-2a\chi}\bigl(\Scal_{\widehat h}
       -4a\Delta_{\widehat h}\chi
       -2a^2|\dd\chi|_{\widehat h}^2\bigr).
\end{equation}
Choose $a$ so that
$4a\|\Delta\chi\|_\infty+2a^2\|\dd\chi\|_\infty^2\le1$ and
$4a^2|\partial E_0|\le1$.  Further decreasing $a$ gives any prescribed
$C^2$ or tensor-comparison tolerance.  There is no need for a uniform
collar width or a uniform positive lower bound for the initial area.
This modification will be used only for the flow calculation; mass
inequalities are applied to the original comparison metric.
\end{proof}

\subsection{Weak flow and its topology}

We specify the weak-flow input.  Theorem~3.1 of
\citet{HI2001} applies on a smooth complete connected
boundaryless manifold possessing a proper locally Lipschitz weak
subsolution with precompact initial set.  It supplies a proper weak
solution from every nonempty smooth precompact open set.  In dimension
three, Regularity~1.3, Hull Property~1.4, and Growth Lemma~1.6 give
compact $C^{1,\gamma}$ level boundaries for $0<\gamma\le1/2$,
outward minimality, and
exponential area growth from a minimizing hull.  Smooth Start
Lemma~2.4 applies when that hull is strict and its smooth boundary has
$H>0$.  We use the $H^2$ inequality from Step~1 of the proof of
Growth Formula~5.7, Equation~(5.22): it holds for almost every pair of
positive endpoints and precompact levels, without a connectedness or
scalar-curvature-sign assumption.  This form suffices for a
distributional inequality in time.  Our smooth initial adjustment
and Smooth Start Lemma~2.4 supply its initial upper trace.  Thus no
restart from an arbitrary nonsmooth positive-time level is needed.

The end hypothesis in this existence theorem holds for all metrics
used here.  On an $\AH$ end the coordinate-sphere mean curvature is
$2+o(1)$ and $|\dd\log r|=1+o(1)$; on an $\AF$ end the corresponding
quantities are $2/r+o(r^{-1})$ and $r^{-1}+o(r^{-1})$.
Thus $u_\infty=c\log r$ with $0<c<2$ has
\begin{equation}\label{eq:flow-end-subsolution}
 \divg_{\widehat h}\left(
    \frac{\nabla u_\infty}{|\nabla u_\infty|}\right)
 \ge |\nabla u_\infty|
\end{equation}
beyond a sufficiently large sphere $r=R$.  Use
$v=\max\{c\log(r/R),0\}-1$, extended by $-1$ throughout the inner
region.  Its negative sublevel is a nonempty smooth precompact set,
bounded by $r=R e^{1/c}$, and outside this set the displayed inequality
holds smoothly.  Extending the radial unit normal field by zero
inside $r=R$ gives a nonnegative distributional divergence
contribution at that sphere.  The Gauss--Green inequality against
outward variations proves the weak subsolution property as well.
This subsolution is proper toward infinity.  The compact
metric adjustment in Lemma~\ref{lem:flow-initial-adjustment} leaves it
unchanged.

For positive-time notation, extend the exterior level function $u$ by
zero on the initial set and, for $t>0$, write
\begin{equation}\label{eq:flow-levels}
 E_t=E_0\cup\{u<t\},\qquad N_t=\partial E_t,\qquad
 A(t)=|N_t|_{\widetilde h}=e^t A(0).
\end{equation}
Use the open perimeter representatives of these sets.  At almost every
time their boundary is $C^{1,\gamma}$, its weak mean curvature is
$H=|\nabla u|$, and its area measure agrees with the level measure in
the coarea formula.  Indeed, differentiability and nonvanishing
gradient hold at almost every point of almost every level in the
coarea sense; those points belong to the reduced boundary of the
sublevel set.  The equivalent statement follows directly from the BV
coarea formula.  The variational property used below is local
minimality of
\begin{equation}\label{eq:flow-variational-functional}
 F\longmapsto \Per_{\widetilde h}(F)
       -\int_F|\nabla u|_{\widetilde h}\,\dd V_{\widetilde h}
\end{equation}
under compact changes away from the initial set; differences of this
expression are taken in a compact region containing the change.

\begin{lemma}[A component bound on fixed topology]\label{lem:hull-components}
At almost every positive time, every component of $E_t$ meets $E_0$,
the open exterior of $E_t$ is connected, and
\begin{equation}\label{eq:flow-component-bound}
 b_0(N_t)\le b_2,\qquad \chi(N_t)\le2b_2,
\end{equation}
where $b_2$ is the fixed number in
\eqref{eq:flow-topology-constant}.
\end{lemma}

\begin{proof}
Fix a time $t$ with the stated regularity and with
$\Vol\{u=t\}=0$.  Excluding the latter exceptional times removes at
most a countable set.  Suppose a component $B$ of $E_t$ did not meet
$E_0$.  For almost every $0<s<t$, the set $E_s\cap B$ is compactly
contained in $B$ and separated from $E_0$.  To check this, continuity
gives $u\le s<t$ on its closure, and a neighborhood of
$\overline E_0$ belongs to $E_t$.  A connected component of that
neighborhood cannot meet $B$ without placing part of $E_0$ in $B$.
The boundaries here are perimeter representatives; the same
containment holds modulo null sets.

Removing $E_s\cap B$ is consequently an allowed variation in
\eqref{eq:flow-variational-functional}.  If
$p(s)=\Per_{\widetilde h}(E_s;B)$, minimality and coarea give
\begin{equation}\label{eq:flow-no-new-components}
 p(s)\le\int_{E_s\cap B}|\nabla u|\,\dd V
       =\int_0^s p(\ell)\,\dd\ell.
\end{equation}
There is no initial contribution because $B\cap E_0=\varnothing$.
Gr\"onwall's inequality implies $p=0$ almost everywhere.  A
finite-perimeter indicator with zero variation on the connected open
set $B$ is constant there.  Compact containment rules out the value
one, so $E_s\cap B$ is null for almost every $s<t$.  Letting such $s$
increase to $t$ contradicts that $B$ is a nonempty open component of
$E_t$.  This proves the first assertion.  It also shows that each such
component contains an entire component of $E_0$, and hence meets $O$.

Outward minimality excludes a bounded open component of the
complement: filling it removes its nonempty boundary and strictly
decreases perimeter.  Because $X$ has just one product end and $E_t$
is precompact, only one component of the exterior can be unbounded.
Its open exterior is therefore connected.

Let $N_t^1,\ldots,N_t^m$ be the boundary components.  Choose a large
compact truncation
$Y\subset X\setminus\operatorname{int}O$ containing all of them.
For each $j$, an arc can start on $\partial O$, travel in the adjacent
component of $E_t$, cross $N_t^j$ once transversely, and then travel in
the connected exterior to the outer boundary of $Y$, without meeting
any other $N_t^k$.  To construct its inner part, join a point just
inside $N_t^j$ to $O$ in the relevant connected open component and
stop at the first cap encountered; thus the arc stays outside the cap
interiors.  Small perturbations give the asserted transversality.

These arcs define relative one-cycles in $(Y,\partial Y)$ whose
mod-two intersection pairings with the closed surfaces $N_t^k$ are
the separate coordinate vectors.  Hence the classes $[N_t^k]$ are
linearly independent in $H_2(Y;\mathbb F_2)$.  Extending the
truncation along its product end changes no homology, so $m\le b_2$.
Each closed orientable surface has Euler characteristic at most two,
which proves the second estimate.  The arcs and the intersection
pairings also make explicit why disconnected caps and arbitrary
compact interior topology cause no difficulty.
\end{proof}

\begin{lemma}[Curvature control up to a prescribed area]\label{lem:flow-curvature}
Suppose the adjusted metric satisfies
$\Scal_{\widetilde h}\ge-C_0$ outside $E_0$, and let $A_*>0$.
At almost every time for which $A(t)\le A_*$,
\begin{equation}\label{eq:flow-uniform-curvature}
 J(t):=\int_{N_t}|\nabla u|^2\,\dd A_{\widetilde h}
 \le C_H:=1+\frac23 C_0 A_*+16\pi b_2.
\end{equation}
In particular this bound is independent of $A(0)$, even if the
initial areas tend to zero in a family.
\end{lemma}

\begin{proof}
Let $\mathcal A_t$ denote the weak second fundamental form of $N_t$.
For almost every pair $0<r<t$, the cited growth formula, after
dropping its nonnegative tangential gradient term, says
\begin{equation}\label{eq:flow-hi-growth}
 J(r)\ge J(t)+\int_r^t\int_{N_s}
 \bigl(2|\mathcal A_s|^2+2\Ric(\nu,\nu)-H^2\bigr)
 \,\dd A\,\dd s.
\end{equation}
After our adjustment, Smooth Start Lemma~2.4 identifies the weak
flow with the smooth flow on an initial interval.  Hence
$J(r)\to J(0)\le1$ as $r\downarrow0$.  Letting admissible
Lebesgue endpoints tend to zero supplies the initial upper trace in
the integrated inequality below.  At positive times only the
almost-everywhere inequality is used.

For completeness, Gauss and Gauss--Bonnet apply to these levels with
their weak curvature.  A local $C^{1,\gamma}$ graph with bounded weak
mean curvature belongs to $W^{2,2}$ by difference quotients in its
uniformly elliptic mean-curvature equation.  Smoothing graph
representations gives smooth embeddings converging in
$C^1\cap W^{2,2}$; the graph formula for the second fundamental form
then gives convergence of the quadratic curvature terms in $L^1$.
Thus the smooth Gauss identity and Gauss--Bonnet pass to the limit.
This is also the content of Lemmas~5.4--5.5 of
\citet{HI2001}.

The Gauss equation and $|\mathcal A_s|^2\ge H^2/2$ give
\begin{align}\label{eq:flow-gauss-bound}
 \int_{N_s}\bigl(2|\mathcal A_s|^2+2\Ric(\nu,\nu)-H^2\bigr)
 &=\int_{N_s}\bigl(|\mathcal A_s|^2+
                   \Scal_{\widetilde h}-\Scal_{N_s}\bigr)\notag\\
 &\ge\tfrac12J(s)-C_0A(s)-8\pi b_2.
\end{align}
Consequently, in distributions in time,
$J'+J/2\le C_0A+8\pi b_2$, with initial upper value at most one.
One can see this directly from \eqref{eq:flow-hi-growth} at Lebesgue
endpoints: $J$ plus the integral of the lower bound in
\eqref{eq:flow-gauss-bound} has a nonincreasing representative.  Any
singular part of its derivative is nonpositive, so the integrating
factor is valid.  It yields
\begin{equation}\label{eq:flow-integrated-curvature}
 J(t)\le e^{-t/2}
       +\int_0^t e^{(s-t)/2}\bigl(C_0A(s)+8\pi b_2\bigr)\,\dd s.
\end{equation}
Using $A(s)=e^{s-t}A(t)$ bounds the two integrals by
$\frac23C_0A(t)$ and $16\pi b_2$, respectively.
\end{proof}

\subsection{Projecting enclosing surfaces}

In the next lemma, $A_*$ is the original minimum enclosing area,
defined using complete intrinsic boundaries.  The comparison end may
be the original $\AH$ end, or the $\AF$ extension of a fixed compact
inner region.  In the latter case $g_{\mathrm{base}}$ denotes $g$ on
the inner region and $g_o$ on the extension.  Statements involving
this base metric are interpreted separately on the two sides of the
joining sphere $\mathcal C$.

\begin{lemma}[Enclosing-area projection]\label{lem:area-projection}
Let $N$ be a compact embedded $C^1$ enclosure of $O$, disjoint from
$S$, with connected exterior.
\begin{enumerate}[label=\textup{(\roman*)}]
\item On the original $\AH$ exterior, fix a large radius $r_p$ and put
\begin{equation}\label{eq:flow-ah-weight}
 \rho=\min\{1,r_p^2/r^2\},
\end{equation}
with $\rho=1$ on the compact part.  There is $q(r_p)\downarrow0$,
depending only on the original end metric, such that
\begin{equation}\label{eq:flow-ah-projection}
 \int_N\rho\,\dd A_g\ge(1-q(r_p))A_*.
\end{equation}
\item Suppose the $\AF$ extension has the parallel-surface form
$g_o=U^2\dd t^2+\gamma_t$, where $\gamma_t$ is induced on outward
Euclidean parallels of a strictly convex embedding of
$(\mathcal C,\gamma_0)$.  Then, with $\rho=1$,
\begin{equation}\label{eq:flow-af-projection}
 \int_N\dd A_{g_{\mathrm{base}}}\ge A_*.
\end{equation}
For a surface crossing the joining sphere, the assertion holds for
the piecewise area; it is enough that its intersection with the
joining sphere have zero two-dimensional measure.  In particular
this holds for almost every weak-flow level.
\end{enumerate}
\end{lemma}

\begin{proof}
First take $N$ smooth and transverse to the spheres used in the
construction.  Let $D$ be its closed end-side domain.  For part
\textup{(i)}, choose a transverse sphere $r=r_1$ with
$r_p/2\le r_1\le r_p$, and replace $D$ by
\begin{equation}\label{eq:flow-projected-domain}
 D'=D\cup\{r\ge r_1\}.
\end{equation}
This domain is connected because both sets contain the entire remote
end.  Its boundary inside $r_1$ is part of $N$; the remaining boundary
is a portion of the sphere $r=r_1$.  Every point of this newly exposed
sphere portion is covered by radial projection of $N\cap\{r\ge r_1\}$:
the ray starts outside $D$ and eventually lies in its end, so must
cross $N$.

In the hyperbolic reference metric the radial map from radius $r$ to
$r_1$ has two-dimensional Jacobian at most $r_1^2/r^2$ on every
two-plane.  The original metric has uniform relative error $o(1)$
on $r\ge r_p/2$.  Its corresponding Jacobian is therefore at most
\begin{equation}\label{eq:flow-projection-jacobian}
 (1+o_{r_p}(1))\frac{r_1^2}{r^2}
 \le(1+o_{r_p}(1))\rho.
\end{equation}
The area formula, allowing multiplicity, bounds the new sphere area
by the weighted area of the discarded portion of $N$.

The boundary of $D'$ has transverse corners along finitely many
smooth curves.  In local coordinates given by defining functions for
the two faces, round each quadrant edge in a shrinking tubular
neighborhood, toward the side enlarging $D'$.  The added area tends
to zero: the modification has vanishing width along a fixed compact
edge, with bounded local slopes.  The rounded end-side domain stays
connected and gives an admissible cut, disjoint from $S$.  Since the
definition of $A_*$ uses its full intrinsic boundary, it follows that
$A_*\le(1+o_{r_p}(1))\int_N\rho\,\dd A_g$.  This proves
\eqref{eq:flow-ah-projection}.

For part \textup{(ii)}, keep the part of $D$ in the original compact
inner region and attach the whole original $\AH$ region outside
$\mathcal C$.  The resulting end-side domain is connected.  Indeed,
every point of its inner part can be joined inside the old $D$ to
infinity, and the initial portion of such a path reaches
$\mathcal C$.  Its newly exposed boundary in $\mathcal C$ is covered
by projection of the old outer boundary along the parallel rays.
This projection has two-dimensional Jacobian at most one:
$\gamma_t\ge\gamma_0$ by strict convexity, its normal derivative is
zero, and the positive coefficient $U^2$ creates no mixed terms.
Rounding the resulting transverse corners as above gives admissible
original cuts with area at most
$\int_N\dd A_{g_{\mathrm{base}}}$ plus a vanishing error.

Finally, a $C^1$ embedded enclosing surface can be approximated by
smooth embeddings in $C^1$, within a collar disjoint from $S$, by an
ambient isotopy.  This preserves the end side and its connectedness.
Choose the approximations transverse to the cutting sphere.  Areas
and weighted areas converge.  In the joined case, first approximate
on the two smooth sides and make the approximations transverse to
$\mathcal C$; round across a shrinking seam neighborhood.  The base
metric is bounded and smooth on either side at this fixed
construction, and the surface has zero area in the seam.  The area
of its portion in that neighborhood tends to zero, so these local
roundings contribute a vanishing area error.  This procedure also
applies when the smooth collar coordinates chosen for the comparison
metric differ from the original side coordinates: the
identifications are smooth on each closed side and preserve the
joining surface.  For almost every flow level the zero-area
intersection follows by applying the level-measure decomposition to
the fixed finite-area seam.  Passing to the limit proves both
assertions.
\end{proof}

\subsection{The uniform transfer criterion}

The following formulation isolates all requirements on the
deformations.  A reference metric may be piecewise smooth across one
fixed compact joining surface, as in
Lemma~\ref{lem:area-projection}\textup{(ii)}.  Its area and volume
measures below are then the sums of the two side measures.

\begin{proposition}[Uniform area transfer]\label{prop:area-transfer}
Fix the capped topology $(X,O)$, a number $A_*>0$, and a scalar lower
bound $C_R\ge0$.  Suppose a comparison has the following properties.
\begin{enumerate}[label=\textup{(\roman*)}]
\item The smooth complete metric $\widehat h$ has one $\AH$ or $\AF$
end with the mean-curvature and radial-gradient asymptotics used in
\eqref{eq:flow-end-subsolution}.  The cap-only enclosure $E_0$ is a
smooth precompact strictly outward-minimizing set, with minimal
boundary disjoint from $O$, every component meeting $O$, and
$\Scal_{\widehat h}\ge-C_R$ outside $E_0$.
\item On $X\setminus O$ there are a reference metric
$g_{\mathrm{base}}$, positive functions $x,W$, a metric $\bar g$, and
a weight $0<\rho\le1$, such that
\begin{equation}\label{eq:flow-metric-comparison}
 \bar g\ge g_{\mathrm{base}},\qquad
 \dd V_{\bar g}=W\,\dd V_{g_{\mathrm{base}}},\qquad
 (1-\eta)x\bar g\le\widehat h\le(1+\eta)x\bar g.
\end{equation}
\item Every enclosing $C^1$ surface with connected exterior, as in
Lemma~\ref{lem:area-projection}, satisfies
\begin{equation}\label{eq:flow-abstract-projection}
 \int_N\rho\,\dd A_{g_{\mathrm{base}}}\ge(1-q)A_*.
\end{equation}
In the piecewise case this is required for surfaces having zero area
in the joining surface.
\end{enumerate}
For every $0<\eps<1$ there are positive tolerances
$\eta_0,q_0,\delta_0$ and $v_0$, depending only on
$\eps,A_*,C_R$ and the integer $b_2$ in
\eqref{eq:flow-topology-constant}, with the following consequence.
If $\eta\le\eta_0$, $q\le q_0$, $0<\delta\le\delta_0$, and
\begin{equation}\label{eq:flow-bad-volume}
 \int_{\{x<1-\delta\}}\rho^{3/2}W\,\dd V_{g_{\mathrm{base}}}<v_0,
\end{equation}
then
\begin{equation}\label{eq:flow-area-conclusion}
 |\partial E_0|_{\widehat h}\ge(1-\eps)A_*.
\end{equation}
The tolerances are uniform over all comparisons satisfying these
conditions; the geometry and initial area of $E_0$ may vary.
\end{proposition}

\begin{proof}
Make the change from Lemma~\ref{lem:flow-initial-adjustment}.  It
preserves the initial area and all topological properties.  Take its
size sufficiently small that the comparison
\eqref{eq:flow-metric-comparison} for the resulting $\widetilde h$
holds with an error $\bar\eta\le\eps/32$, by initially requiring,
for example, $\eta\le\eta_0=\eps/64$.  Its scalar lower bound is
$-C_0$, where $C_0=C_R+1$.  Also take
$q_0=\delta_0=\eps/32$.

Suppose, to the contrary, that $A(0)<(1-\eps)A_*$.  Exponential area
growth gives an interval $I=[t_-,t_+]$ on which the endpoint areas
are $(1-\eps)A_*$ and $(1-\eps/2)A_*$.  Its length is the fixed
positive number
\begin{equation}\label{eq:flow-time-window}
 \Delta_\eps=\log\frac{1-\eps/2}{1-\eps}.
\end{equation}
Although $t_-$ may tend to infinity when $A(0)$ tends to zero,
Lemma~\ref{lem:flow-curvature} bounds $J(t)$ throughout $I$ by the
same constant $C_H$.

Write $\mathcal B=\{x<1-\delta\}$.  On the good part of a level,
\eqref{eq:flow-metric-comparison} and $\bar g\ge g_{\mathrm{base}}$
give
\begin{equation}\label{eq:flow-good-area}
 \int_{N_t\setminus\mathcal B}\rho\,\dd A_{g_{\mathrm{base}}}
 \le\frac{A(t)}{(1-\bar\eta)(1-\delta)}.
\end{equation}
Combine this with \eqref{eq:flow-abstract-projection}.  Since
$(1-\bar\eta)(1-\delta)\ge1-\eps/16$, for almost every $t\in I$
we obtain
\begin{equation}\label{eq:flow-bad-area-lower}
 \int_{N_t\cap\mathcal B}\rho\,\dd A_{g_{\mathrm{base}}}
 \ge (1-q)A_*
       -\frac{(1-\eps/2)A_*}{(1-\bar\eta)(1-\delta)}
 \ge \frac{\eps A_*}{4}=:c_A.
\end{equation}
The same metric comparison on the bad part therefore implies
\begin{equation}\label{eq:flow-weighted-area-lower}
 \int_{N_t\cap\mathcal B}\rho x^{-1}\,\dd A_{\widetilde h}
 \ge(1-\bar\eta)c_A\ge c_A/2.
\end{equation}

Here is the precise weighted coarea step.  Put $F=\rho x^{-1}$ and
$H=|\nabla u|_{\widetilde h}$.  For almost every level $H>0$ at
almost every point in its area measure.  H\"older with exponents
$3/2$ and $3$ gives
\begin{align}\label{eq:flow-holder-exponents}
 \int_{N_t\cap\mathcal B}F\,\dd A
 &=\int_{N_t\cap\mathcal B}
       \left(\frac{F^{3/2}}H\right)^{2/3}(H^2)^{1/3}\,\dd A\notag\\
 &\le
 \left(\int_{N_t\cap\mathcal B}\frac{F^{3/2}}H\,\dd A\right)^{2/3}
 J(t)^{1/3}.
\end{align}
Infinite integrals cause no difficulty; otherwise this identity can
first be applied where $H$ is bounded away from zero and passed to
the limit.  Equations~\eqref{eq:flow-uniform-curvature} and
\eqref{eq:flow-weighted-area-lower} yield
\begin{equation}\label{eq:flow-level-volume-lower}
 \int_{N_t\cap\mathcal B}
       \frac{\rho^{3/2}x^{-3/2}}{|\nabla u|}\,\dd A_{\widetilde h}
 \ge \frac{(c_A/2)^{3/2}}{\sqrt{C_H}}=:c_V.
\end{equation}
Integrating over $I$ and applying Lipschitz coarea gives
\begin{align}\label{eq:flow-coarea-contradiction}
 \Delta_\eps c_V
 &\le\int_{\mathcal B\cap\{t_-<u<t_+\}\cap\{|\nabla u|>0\}}
                  \rho^{3/2}x^{-3/2}\,\dd V_{\widetilde h}\notag\\
 &\le(1+\bar\eta)^{3/2}
       \int_{\mathcal B}\rho^{3/2}W\,\dd V_{g_{\mathrm{base}}}.
\end{align}
The last inequality uses
$\dd V_{\widetilde h}\le
(1+\bar\eta)^{3/2}x^{3/2}W\,\dd V_{g_{\mathrm{base}}}$.
Thus it is enough to take
$v_0=\Delta_\eps c_V/2^{3/2}$.
Notice that no contribution from a plateau of the level function
has been assumed: coarea integrates only where its gradient is
nonzero, and the full bad-set volume is an upper bound.  This proves
the proposition.
\end{proof}

\subsection{The two applications and their parameter order}

For the maximal $\AH$ reduction, the base is the original $g$ and
\begin{equation}\label{eq:flow-graph-volume}
 \bar g=g+s^2\dd f^2,\qquad h=x\bar g,\qquad
 W^2=1+s^2|\dd f|_g^2,\qquad
 \dd V_{\bar g}=W\,\dd V_g.
\end{equation}
The curvature bound is $\Scal_h\ge-6$, and the cap-only largest hull
is the enclosure used in that reduction.  Fix the area tolerance
$\eps$, then $r_p$ so that the error in
\eqref{eq:flow-ah-projection} is at most $q_0$, and then
$\delta\le\delta_0$.  The quantitative estimates of
Section~\ref{sec:maximal-reduction} give, with
$R_s=M^{3/10}$ and $L=M^2$,
\begin{equation}\label{eq:flow-ah-bad-volume}
 \int_{\{x<1-\delta\}}\rho^{3/2}W\,\dd V_g
 \le C\left(L^{-1}+\sqrt{M/L}
              +R_s^{-1}+\frac{\sqrt M}{R_s^2}\right)
 \longrightarrow0.
\end{equation}
To explain the powers, the unweighted bad volume inside $R_s$ is
$O(L^{-1})$, and $\int(W^2-1)\,\dd V_g\le CM$; Cauchy--Schwarz
gives the first two terms.  On the tail the $\AH$ volume element is
comparable to $r\,\dd r\,\dd A_\sigma$ and
$\rho^{3/2}=r_p^3/r^3$.  Hence
$\int_{r>R_s}\rho^{3/2}\,\dd V_g=O(R_s^{-1})$, whereas
$\int_{r>R_s}\rho^3\,\dd V_g=O(R_s^{-4})$; another
Cauchy--Schwarz inequality gives the last term.  The constants here
may depend on the already fixed $r_p$ and original data.
Proposition~\ref{prop:area-transfer} therefore yields the required
lower bound for the comparison boundary area.

For the compact construction with the $\AF$ attachment,
$g_{\mathrm{base}}=g$ inside and $g_o$ outside,
$\bar g=g+s^2\dd f^2$ inside and $\bar g=g_o$ outside.  Thus $W=1$
outside, $\rho=1$, and Lemma~\ref{lem:area-projection} has zero
projection error.  Section~\ref{sec:scalar-error} constructs a smooth
metric $h'$ arbitrarily close to $h=x\bar g$, with a uniform scalar
lower bound, together with two enclosures.  The enclosure $E$
contains the caps and all unsafe balls and is used for the
$\AF$ mass inequality.  The largest cap-only minimizer $E_0$ is
used for the flow.  In particular
\begin{equation}\label{eq:flow-af-two-enclosures}
 |\partial E|_{h'}\ge|\partial E_0|_{h'}.
\end{equation}
No topological bound on the unsafe balls or on $E$ is required.

The bad set is confined to a fixed outer radius by the capacitary
lower bound for $x$.  On its designated penalized region,
\begin{equation}\label{eq:flow-af-bad-volume}
 \int_{\{x<1-\delta\}\cap\mathrm{designated}}
                       W\,\dd V_{g_{\mathrm{base}}}\le C/L.
\end{equation}
The omitted fixed exterior collar has $W=1$, so its contribution is
at most its base volume.  The omitted inner collar has base width
$O(M^{-4})$.  Since the compact construction gives
$\int_{\Omega_i}W^2\,\dd V_g\le C(L)M^2$, its contribution is at
most
\begin{equation}\label{eq:flow-inner-collar-cost}
 O(M^{-2})\left(\int_{\Omega_i}W^2\,\dd V_g\right)^{1/2}
 =O_L(M^{-1}).
\end{equation}
Choose the exterior collar sufficiently thin and then $L$
sufficiently large, both before $M$.  Next choose $M$ large and the
remaining finite penalty parameters as in
Section~\ref{sec:compact-construction}.  Finally choose the individual
corner-smoothing tolerance small enough for
\eqref{eq:flow-metric-comparison}.  These choices make the bad volume
smaller than $v_0$.  The flow-only initial adjustment does not change
$|\partial E_0|_{h'}$.  Equations~\eqref{eq:flow-area-conclusion} and
\eqref{eq:flow-af-two-enclosures} give
\begin{equation}\label{eq:flow-af-area-conclusion}
 |\partial E|_{h'}\ge|\partial E_0|_{h'}\ge(1-\eps)A_*.
\end{equation}
If $A_*=0$ there is no area-transfer assertion to prove; the
nonnegative mass conclusion from the $\AF$ comparison supplies the
needed zero-area bound.
